\section{The cost model and the crossover payload}
\label{sec:costmodel}

\begin{figure*}[t]
\centering
\lrfigure{cost_model}{\linewidth}
\caption{Cycles per packet against payload size, with the fitted
$C(s)=a+b\,s$. Bars are bootstrap intervals on the median of \envReplicates\
replicates; hollow markers are points the stability gate excluded from the fit.
Left, the hardware crypto instructions available; right, the same measurement
with them masked in both implementations.}
\label{fig:costmodel}
\end{figure*}

\begin{table*}[t]
\caption{Fitted cost model. $a$ is the payload-independent cost, $b$ the cost
per payload byte over a path that crosses the cipher twice, and $s^{*}=a/b$ the
crossover payload with its 95\% interval propagated through the parameter
covariance.}
\label{tab:costmodel}
\centering\footnotesize
\input{generated/table_cost_model}
\end{table*}

The sweep comprises \eOneUnits\ distinct configurations.
Figure~\ref{fig:costmodel} shows it and Table~\ref{tab:costmodel} the fits. The
relationship is linear over the whole payload range for every cell but one. For
\bestCipher\ with the hardware instructions masked, \fitCurvatureLrAesGcmMasked.

The fits are weighted by the inverse variance of each point rather than left
unweighted. Small payloads are noisier, because that is where $a$ dominates and
where scheduling noise lands, and an unweighted fit would let the noisiest
points set the intercept, which is the parameter the whole study turns on. The
interval on $s^{*}$ is propagated through the covariance of $a$ and $b$ rather
than obtained by dividing their separate intervals. The two are strongly
anticorrelated, since a fit that raises the intercept lowers the slope, and
ignoring that inflates the interval several-fold.

\begin{figure*}[t]
\centering
\lrfigure{crossover}{\linewidth}
\caption{The crossover payload for each cell of the design, with its 95\%
interval. The dashed line is the largest payload this data plane can carry
inside an MTU: a crossover to the right of it means the fixed cost dominates at
every size the data plane can send.}
\label{fig:crossover}
\end{figure*}

\noindent\textbf{The headline.}
With hardware crypto available, \bestCipher\ has $a = \bestA$ cycles and
$b = \bestB$ cycles per payload byte, giving $s^{*} = \bestSStar$ bytes. The
largest payload this data plane can carry inside an MTU is \protoMaxPayload\
bytes, so the crossover lies outside the reachable range by a factor of more
than two. The fixed cost therefore dominates at every payload size this
configuration can send, and no improvement to the cipher can change its
throughput by much: the packet costs, not the byte. Where $s^{*}$ falls inside
the range instead, as it does for \texttt{lr-chacha20-poly1305} in
Table~\ref{tab:costmodel}, the crossover is reachable and the choice of cipher
starts to matter above it.

\noindent\textbf{Reading $a$ across experiments.}
Table~\ref{tab:costmodel} obtains $a$ from a fit across the whole payload sweep,
while Table~\ref{tab:controls} and Table~\ref{tab:virtualization} report a cost
measured at one configuration. The three are separate experiments on the same
host and their values for the same cipher agree to about one per cent, which is
the reproducibility of the harness rather than a discrepancy. Comparisons in
this report are always drawn within one experiment.

\noindent\textbf{Our implementations against OpenSSL's.}
The AES-256-GCM written for this study reaches $b$ a factor of \simdVsOpensslB\
of OpenSSL's, with a fixed cost a factor of \simdVsOpensslA\ of it. Where that
factor exceeds one it has a specific cause, and we state it rather than
attributing it to implementation quality in general. OpenSSL's assembly
interleaves eight AES blocks using the vector AES extensions, while the
implementation here interleaves four using the 128-bit instructions. The
stitched inner loop, which computes counter-mode encryption and the GHASH in one
pass over the data rather than two, recovers most but not all of that
difference. Our implementation is checked byte for byte against OpenSSL at every
length in the sweep, so the residual is a property of the instruction selection
and not of correctness.

The comparison is useful in both directions. It bounds how much of the measured
$b$ is intrinsic to the algorithm rather than to one library, and it puts a
number on what the wider vector path buys, which is invisible when only one
implementation is measured.

\noindent\textbf{Masking the accelerator.}
The masked arm is not a hypothetical. It is what the same code does on a machine
without the instructions, which is what an embedded or older deployment target
looks like. The change in $b$ between the two arms is the whole contribution of
the hardware accelerator to the bulk term, and Table~\ref{tab:costmodel} shows
it spanning more than an order of magnitude. The change in $a$ is smaller. It is
the accelerator's contribution to the cipher's per-message setup, the part that
does not scale with payload and therefore does not go away by sending larger
packets.
